\documentclass[11pt]{article}
\usepackage{amsmath,amssymb,amsthm}
\usepackage[margin=1in]{geometry}

\newtheorem{definition}{Definition}
\newtheorem{proposition}{Proposition}
\newtheorem{remark}{Remark}

\title{Hyperseed Formalization:\\Sketch of an AGI Curriculum}
\author{ProtomegaTron (automated interpretation)}
\date{2026-09-18}

\begin{document}
\maketitle

\section{Overview}

This note formalizes the structural content of Goertzel's 2021 AGI
curriculum sketch within the Hyperseed ontology. The key insight is that
an AGI curriculum is itself a fiber bundle atlas over the knowledge base
space, and that common mistakes in AI correspond to specific
fiber-collapse errors.

\section{Curriculum as Base-Space Atlas}

\begin{definition}[AGI Knowledge Base Space]
Let $B_{\text{AGI}}$ be the base space of AGI-relevant knowledge. The
curriculum partitions $B_{\text{AGI}}$ into six overlapping charts:
\[
  B_{\text{AGI}} = U_{\text{hist}} \cup U_{\text{alg}} \cup
  U_{\text{neuro}} \cup U_{\text{phil}} \cup U_{\text{arch}} \cup
  U_{\text{future}}
\]
where each $U_i$ corresponds to one of the six curriculum areas.
\end{definition}

\begin{definition}[Cross-Cutting as Transition Functions]
Books that span multiple categories provide transition functions on
chart overlaps:
\[
  \phi_{ij}: U_i \cap U_j \to GL(F)
\]
where $F$ is the fiber of understanding and $GL(F)$ is the structure
group. The ``crude division'' Goertzel notes reflects the fact that
$U_i \cap U_j \neq \emptyset$ for many pairs $(i,j)$.
\end{definition}

\section{Mitchell's Mistakes as Fiber-Collapse Errors}

\begin{proposition}[Four Fiber-Collapse Errors]
Mitchell's four recurring AI mistakes correspond to specific structural
errors in the fiber bundle:
\begin{enumerate}
  \item \textbf{Narrow/general continuum:} Collapsing distinct fibers
    $F_{\text{narrow}}$ and $F_{\text{general}}$ into one, losing the
    structural distinction between narrow competence and general
    understanding.
  \item \textbf{``Easy things are easy'':} Projecting base-space
    difficulty $d_B(x,y)$ without checking the fiber structure ---
    tasks easy in the base may require complex fiber transport, and
    vice versa.
  \item \textbf{Wishful mnemonics:} Labeling a fiber as ``perception''
    or ``understanding'' without verifying that the labeled process
    actually implements the corresponding transport map.
  \item \textbf{Brain-only intelligence:} Restricting the base space to
    $B_{\text{brain}} \subset B_{\text{AGI}}$, losing the embodied,
    relational, and social fibers that contribute to general
    intelligence.
\end{enumerate}
\end{proposition}

\section{d-Calculus: Curriculum as d-Connected Space}

\begin{definition}[Disciplinary Distinction Metric]
Define the distinction metric $d: B_{\text{AGI}} \times B_{\text{AGI}}
\to \mathbb{R}_{\geq 0}$ where $d(x,y)$ measures the conceptual
distance between topics $x$ and $y$. The curriculum is well-designed
if $(B_{\text{AGI}}, d)$ is connected: for any two topics $x, y$,
there exists a path $x = z_0, z_1, \ldots, z_n = y$ with
$d(z_i, z_{i+1}) < \epsilon$ for all $i$, where each step corresponds
to a reading or assignment in the curriculum.
\end{definition}

\begin{proposition}[d-Distortion from Mistakes]
Each of Mitchell's mistakes induces a specific d-distortion:
\begin{itemize}
  \item Narrow/general continuum: $d(F_{\text{narrow}}, F_{\text{general}})
    \to 0$ (collapse of genuine distinction).
  \item ``Easy things are easy'': $d_{\text{perceived}}(x,y) \neq
    d_{\text{actual}}(x,y)$ (systematic projection error).
  \item Wishful mnemonics: $d(\text{label}, \text{reality}) > 0$ but
    treated as $0$ (mislabeling).
  \item Brain-only: $d(x, B_{\text{embodied}}) \to \infty$ for
    $x \in B_{\text{brain}}$ (artificial restriction).
\end{itemize}
\end{proposition}

\section{Hyperon as Composed Bundle}

\begin{remark}
The curriculum's coverage of evolutionary, neural, logical, and
probabilistic methods mirrors the Hyperon architecture's composition of
corresponding fiber bundles. Goertzel's claim that ``OpenCog Hyperon
embedded in SingularityNET goes beyond'' Mitchell's mistakes is the
claim that the composed bundle $E_{\text{Hyperon}}$ avoids all four
fiber-collapse errors by maintaining distinct fibers, checking
transport, verifying labels, and operating on the full
(embodied/social/distributed) base space.
\end{remark}

\end{document}
